\section*{अध्याय \ref{ch:g10:algebra} --- बीजगणित: समीकरण और असमिकाएँ}

\begin{solution}{exo:g10:algebra:1}
$3(2x-5) - 2(x-7) = 6x - 15 - 2x + 14 = 4x - 1$।

$(x+4)(2x-1) = 2x^2 - x + 8x - 4 = 2x^2 + 7x - 4$।

$(3x-2)^2 = 9x^2 - 12x + 4$।

$(5-x)(5+x) = 25 - x^2$।
\end{solution}

\begin{solution}{exo:g10:algebra:2}
$x^2 - 49 = (x+7)(x-7)$ (वर्गों का अंतर)।

$x^2 + 10x + 25 = (x+5)^2$।

$7x^2 - 14x = 7x(x - 2)$ (उभयनिष्ठ गुणनखंड)।

$16x^2 - 8x + 1 = (4x - 1)^2$।
\end{solution}

\begin{solution}{exo:g10:algebra:3}
$4x + 9 = x - 3$: $3x = -12$, इसलिए $x = -4$।

$\frac{2x-1}{3} = 5$: $2x - 1 = 15$, इसलिए $x = 8$।

$2(x-3) = 2x + 1$ का \omterm{def:g10:algebra:expand}{प्रसार} $2x - 6 = 2x + 1$ है, अर्थात् $-6 = 1$: हर $x$ के लिए असत्य। कोई हल नहीं ---
दोनों पक्ष ऐसी समांतर रेखाएँ बताते हैं जो कभी नहीं मिलतीं।
\end{solution}

\begin{solution}{exo:g10:algebra:4}
$(x+5)(2x-8) = 0$: $x = -5$ या $x = 4$।

$x^2 - 16 = (x-4)(x+4) = 0$: $x = 4$ या $x = -4$।

$x^2 + 6x + 9 = (x+3)^2 = 0$: अकेला हल $x = -3$।
\end{solution}

\begin{solution}{exo:g10:algebra:5}
$3x - 5 \leq 7$: $3x \leq 12$, इसलिए $x \leq 4$: हल-समुच्चय $\intoc{-\infty}{4}$।

$-4x + 1 > 9$: $-4x > 8$, और $-4$ से भाग देने पर असमिका पलट जाती है: $x < -2$: हल-समुच्चय
$\intoo{-\infty}{-2}$।

$2(x-1) \geq 5x + 4$: $2x - 2 \geq 5x + 4$, इसलिए $-3x \geq 6$, इसलिए $x \leq -2$: हल-समुच्चय $\intoc{-\infty}{-2}$।
\end{solution}

\begin{solution}{exo:g10:algebra:6}
उभयनिष्ठ गुणनखंड $(x-1)$ है:
\[
E = (x-1)\bigl[(x+3) - (2x - 5)\bigr] = (x-1)(x + 3 - 2x + 5)
= (x-1)(8 - x).
\]
तब $E = 0$ तब होता है जब $x = 1$ या $x = 8$ हो (शून्य-गुणनफल नियम)।
\end{solution}

\begin{solution}{exo:g10:algebra:7}
$x^2 = 5x$: $x^2 - 5x = x(x - 5) = 0$, इसलिए $x = 0$ या $x = 5$ ($x$ से कभी भाग मत दीजिए: हल $0$ खो जाता)।

$(2x+1)^2 = (x-4)^2$: बाईं ओर लाइए और वर्गों के अंतर का \omterm{def:g10:algebra:expand}{गुणनखंडन} कीजिए:
\[
(2x+1)^2 - (x-4)^2
= \bigl[(2x+1) + (x-4)\bigr]\bigl[(2x+1) - (x-4)\bigr]
= (3x - 3)(x + 5).
\]
शून्य-गुणनफल नियम: $x = 1$ या $x = -5$।
\end{solution}

\begin{solution}{exo:g10:algebra:8}
$(x+2)(4-x) > 0$: गुणनखंड $x + 2$ $-2$ पर लुप्त होता है (उसके बाद धनात्मक), और $4 - x$ $4$
पर लुप्त होता है (उससे पहले धनात्मक)। गुणनफल ठीक तब धनात्मक है जब दोनों
गुणनखंड धनात्मक हों, अर्थात् $\intoo{-2}{4}$ पर।

$\frac{2x-6}{x+1} \geq 0$: अंश $3$ पर शून्य, हर $-1$ पर शून्य (वर्जित)। \omterm{def:g10:algebra:signtable}{चिह्न-सारणी}: भागफल तब
धनात्मक है जब दोनों भागों के चिह्न समान हों, अर्थात् $x < -1$ या $x > 3$ के लिए;
वह $x = 3$ पर लुप्त होता है। हल-समुच्चय: $\intoo{-\infty}{-1} \cup \intco{3}{+\infty}$।
\end{solution}

\begin{solution}{exo:g10:algebra:9}
$x$ महीनों के बाद पहली सेवा का ख़र्च $9x + 15$ है और दूसरी का $12x$। पहली तब
सस्ती है जब
\[
9x + 15 < 12x
\ \Longleftrightarrow\
15 < 3x
\ \Longleftrightarrow\
x > 5 .
\]
$6$-वें महीने से आगे कुल मिलाकर पहली सेवा ही सस्ता विकल्प है।
\end{solution}

\begin{solution}{exo:g10:algebra:10}
वर्जित मान: $x = 1$। $x \neq 1$ के लिए दोनों पक्षों को $x - 1$ से गुणा कीजिए:
\[
x + 3 = 2(x - 1) = 2x - 2
\ \Longleftrightarrow\
x = 5 .
\]
चूँकि $5 \neq 1$ है, हल $x = 5$ है। जाँच: $\frac{5+3}{5-1} = \frac84 = 2$।
\end{solution}

\begin{solution}{exo:g10:algebra:11}
$(n+1)^2 - n^2 = n^2 + 2n + 1 - n^2 = 2n + 1$, जो हर \omterm{def:g10:numbers:sets}{पूर्णांक} $n$ के लिए विषम है।

दो क्रमागत विषम संख्याओं को $2n - 1$ और $2n + 1$ लिखा जा सकता है। तब
\[
(2n+1)^2 - (2n-1)^2
= \bigl[(2n+1)+(2n-1)\bigr]\bigl[(2n+1)-(2n-1)\bigr]
= 4n \times 2 = 8n,
\]
जो $8$ का गुणज है।
\end{solution}

\begin{solution}{pb:g10:algebra:1}
\textbf{1.} $x^2 - 10x + 25 = (x - 5)^2$;
$9x^2 - 49 = (3x - 7)(3x + 7)$;
$x^2 + 6x + 9 - 4 = (x + 3)^2 - 2^2 = (x + 1)(x + 5)$।

\textbf{2.} $x = -1$ या $x = -5$; \quad $x = \frac73$ या $x = -\frac73$।

\textbf{3.} $(5 - t)(5 + t) = 25 - t^2 = 21$ से $t^2 = 4$, $t = 2$ मिलता है: संख्याएँ $3$ और $7$ हैं
(योग $10$, गुणनफल $21$: जाँच लीजिए)।

\textbf{4.} $(7 - t)(7 + t) = 49 - t^2 = 33$: $t = 4$, संख्याएँ $3$ और $11$।

\textbf{5.} $\left(\frac s2 - t\right)\left(\frac s2 + t\right)
= \left(\frac s2\right)^2 - t^2 = p$ से $t^2 = \left(\frac s2\right)^2 - p$ निकलता है। हल ठीक तब है जब यह $\geq 0$ हो, अर्थात्
जब $p \leq \left(\frac s2\right)^2$ हो --- दिए गए योग वाली दो संख्याओं का गुणनफल आधे योग के वर्ग को
कभी नहीं हरा सकता (माध्यमिक विद्यालय खंड की रानी डिडो की बाड़ वाली समस्या
यह जानती थी)।

\textbf{6.} $(x - a)(x - b) = x^2 - (a + b)x + ab$: \omterm{def:g10:algebra:equation}{समीकरण} $x^2 - sx + p = 0$ ठीक यही कहता है कि ``$x$ उन दो संख्याओं
में से एक है जिनका योग $s$ और गुणनफल $p$ है''।

\textbf{7.} $(x-5)^2 - 4 = x^2 - 10x + 25 - 4 = x^2 - 10x +
21$। वर्गों का अंतर: $(x - 5 - 2)(x - 5 + 2) = (x - 7)(x - 3) = 0$: हल $3$ और $7$ --- प्रश्न 3
के लिपिक वाली संख्याएँ।

\textbf{8.} $x^2 + 6x - 7 = (x + 3)^2 - 16 = 0$ से $(x + 3)^2 = 16$, $x + 3 = \pm 4$ मिलता है: $x = 1$ या $x = -7$।

\textbf{9.} हर $x$ के लिए $x^2 + 4x + 5 = (x + 2)^2 + 1 \geq 1 > 0$: कोई वास्तविक हल नहीं। $(x + h)^2 + k$ के रूप में:
हल ठीक तब हैं जब $k \leq 0$ हो (तब $(x+h)^2 = -k$ के दो मूल $-h \pm \sqrt{-k}$ हैं); और जब $k > 0$ हो
तब व्यंजक धनात्मक ही रहता है।

\textbf{10.} भुजा $x$ वाला वर्ग (क्षेत्रफल $x^2$) और दो $3 \times x$ आयत (कुल
$6x$) मिलकर $x^2 + 6x = 7$ क्षेत्रफल का एक L बनाते हैं; छूटा हुआ $3 \times 3$ कोना
(क्षेत्रफल $9$) उसे भुजा $x + 3$ वाला बड़ा वर्ग पूरा कर देता है। दोनों
पक्षों में $9$ जोड़ने पर $(x + 3)^2 = 16$, इसलिए $x + 3 = 4$ (लंबाइयाँ धनात्मक हैं) और
$x = 1$। बीजगणित और चित्र एक ही काम हैं।

\textbf{11.} $(x-3)(x+2)$: मूलों के बीच ऋणात्मक। $< 0$ का हल: $x \in \intoo{-2}{3}$।

\textbf{12.} $x^2 - 10x + 21 \leq 0$, अर्थात् $(x - 3)(x - 7) \leq 0$: मूलों के बीच, सिरे सम्मिलित: $x \in \intcc{3}{7}$।

\textbf{13.} $x = 1$ पर शून्य, $x = -2$ पर वर्जित; भागफल मूलों के बाहर धनात्मक
है: $x \in \intoo{-\infty}{-2} \cup \intco{1}{+\infty}$।

\textbf{14.} $25 - (x - 5)^2 = 25 - x^2 + 10x - 25 =
x(10 - x)$: सर्वसमिका सत्यापित। चूँकि सदा $(x-5)^2 \geq 0$ है, इसलिए $A(x) \leq 25$, और
समता ठीक तब जब $x = 5$ हो: $5 \times 5$ वाला वर्गाकार बाड़ा, अधिकतम $25$ m$^2$ ---
कक्षा 6 की सारणी, अब दो पंक्तियों की उपपत्ति।

\textbf{15.} $x > 0$ के लिए दावे को $x$ से गुणा कीजिए: $x^2 + 1 \geq 2x$, अर्थात् $x^2 - 2x + 1 = (x - 1)^2 \geq 0$
--- जो सत्य है, और समता केवल $x = 1$ पर। वापस $x > 0$ से भाग देने पर सब कुछ बना
रहता है (\cref{prop:g10:algebra:ineqrules})।

\textbf{16.} $\frac1x = \frac{x}{1 - x}$ का तिर्यक गुणन करने पर (सभी लंबाइयाँ धनात्मक हैं):
$x^2 = 1 - x$, अर्थात् $x^2 + x - 1 = 0$।

\textbf{17.} $x^2 + x - 1 = \left(x + \frac12\right)^2 -
\frac54 = 0$ से $x + \frac12 = \pm\frac{\sqrt5}{2}$ मिलता है। धनात्मक मूल: $x = \frac{\sqrt5 - 1}{2} \approx 0.618$।

\textbf{18.} $\varphi = \frac{2}{\sqrt5 - 1} =
\frac{2(\sqrt5 + 1)}{(\sqrt5 - 1)(\sqrt5 + 1)}
= \frac{2(\sqrt5 + 1)}{5 - 1} = \frac{\sqrt5 + 1}{2}
\approx 1.618$।

\textbf{19.} $\varphi^2 = \frac{(\sqrt5 + 1)^2}{4} =
\frac{6 + 2\sqrt5}{4} = \frac{3 + \sqrt5}{2}$, और $\varphi + 1 = \frac{\sqrt5 + 3}{2}$: बराबर। और $\varphi - 1 = \frac{\sqrt5 - 1}{2} = x = \frac1\varphi$ (इसी तरह $\varphi$ परिभाषित किया
गया था)। स्वर्ण अनुपात वह संख्या है जिसका वर्ग एक जोड़ देता है और जिसका
व्युत्क्रम एक घटा देता है।

\textbf{20.} $\sqrt5$ \omterm{ex:g10:numbers:classify}{अपरिमेय} है (माध्यमिक विद्यालय खंड की अपरिमेयता वाली
सप्ताहांत समस्या: $5$ कोई पूर्ण वर्ग नहीं है), और $1$ जोड़कर फिर आधा
करने से अपरिमेयता बनी रहती है (परिमेय संक्रियाएँ, \cref{pb:g10:numbers:1})। अनुपात: $\frac53 \approx 1.667$,
$\frac85 = 1.600$, $\frac{13}{8} = 1.625$, $\frac{21}{13} \approx 1.615$: बारी-बारी से ऊपर और नीचे, $\varphi = 1.618\ldots$ की ओर सिमटते हुए ---
माध्यमिक विद्यालय खंड की लय-गणना वाली सप्ताहांत समस्या की संख्याएँ चुपके से
स्वर्ण विभाजन का पीछा करती हैं; अनुक्रमों वाले अध्याय उन्हें रँगे हाथों
पकड़ेंगे।
\end{solution}
